\par\medskip
\hypertarget{native-079-article060}{}\pdfbookmark[2]{60.}{article-060}
\begin{center}60.\end{center}
Es seien $x$, $y$, $z$ die Abstände des Himmelskörpers von drei Ebenen, die sich in der Sonne rechtwinklig schneiden, und $X$, $Y$, $Z$ die Abstände der Erde (ihres Mittelpunkts oder eines Punktes auf ihrer Oberfläche) von denselben Ebenen. Dann sind $x-X$, $y-Y$, $z-Z$ offenbar die Abstände des Himmelskörpers von drei Ebenen, die parallel zu jenen Ebenen durch die Erde gelegt sind. Diese Abstände stehen zum Abstand des Körpers von der Erde und zu seinem \textit{geozentrischen Ort}\footnote{Im weiteren Sinne. Im eigentlichen Sinne bezieht sich dieser Ausdruck auf den Fall, dass die Gerade vom \textit{Mittelpunkt} der Erde ausgeht.}, also zur Lage der Projektion der Geraden von der Erde zum Körper auf die Himmelskugel, in derselben Beziehung wie $x$, $y$, $z$ zu seinem Abstand von der Sonne und zu seinem heliozentrischen Ort. Es sei $\Delta$ der Abstand des Himmelskörpers von der Erde. Man denke sich auf der \SourceLeaf{080}Himmelskugel ein Lot vom geozentrischen Ort auf den Großkreis gefällt, der der Ebene für die Abstände $z$ entspricht. Es sei $a$ der Abstand des Schnittpunkts vom positiven Pol des Großkreises, der der Ebene für die Abstände $x$ entspricht. Schließlich sei $b$ die Länge des Lots, also der Abstand des geozentrischen Ortes vom Großkreis für die Abstände $z$. Dann ist $b$ die geozentrische Breite beziehungsweise Deklination, je nachdem, ob die Ebene für die Abstände $z$ die Ekliptik oder der Äquator ist. Entsprechend ist $a+N$ die geozentrische Länge beziehungsweise Rektaszension, wenn $N$ im ersten Fall die Länge und im zweiten die Rektaszension des Pols der Ebene für die Abstände $x$ bezeichnet. Somit gilt
\[
\begin{aligned}
x-X&=\Delta\cos b\cos a\\
y-Y&=\Delta\cos b\sin a\\
z-Z&=\Delta\sin b.
\end{aligned}
\]
Die ersten beiden Gleichungen liefern $a$ und $\Delta\cos b$. Die letztere Größe (die positiv anzusetzen ist) liefert zusammen mit der dritten Gleichung $b$ und $\Delta$.
